\setcounter{section}{9}
\section{Datos necesarios para cerrar ingeniería}
\begingroup\small\setlength{\LTpre}{4pt}\setlength{\LTpost}{6pt}
\begin{longtable}{*{3}{L{\dimexpr\linewidth/3-2\tabcolsep\relax}}}
\toprule
\textbf{Prioridad} & \textbf{Dato verificable} & \textbf{Decisión afectada} \\ \midrule
\endfirsthead
\toprule
\textbf{Prioridad} & \textbf{Dato verificable} & \textbf{Decisión afectada} \\ \midrule
\endhead
1 & Placas generador/T1/T2, ubicación y predios & Tensión nominal, potencia neta y transformación \\
1 & Aforos/curva de disponibilidad y demanda horaria & Suficiencia150/200 kW, reserva y horizonte \\
1 & Topografía de campo de terna y cruces & Ubicación, altura y distancia real entre postes \\
1 & Ficha de poste y armadoCAC, cargasN-M & Soportes, retenidas y accesorios \\
1 & Geotecnia y Wenner & Cimentaciones, anclajes y tensiones de paso/contacto \\
1 & AVR/decremento, TC/TT y curvas de equipos & Coordinación de protecciones e inrush \\
2 & Cotizaciones con logística y ensayos & Elección10/22.9 kV y presupuesto \\
2 & Consigna/rúbrica y docente de esta asignatura & Formato de entrega y alcance de evaluación \\
\bottomrule\end{longtable}\endgroup

La ejecución informática y la revisión documental de esta entrega no equivalen a aprobación profesional ni a levantamiento. No se requiere permiso del usuario para continuar corrigiendo cálculos; sí se requieren datos técnicos reales para cerrar las decisiones dependientes. El archivoCONTINUIDAD\_\allowbreak{}IA.md indica cómo reproducir y qué datos actualizar sin reactivar fuentes históricas.

